\chapter{Classification Table}
\label{cha:classification_table}

This table summarizes the 12-dimensional classification system applied to all 841 disease chapters. Each disease is classified across 12 standardized dimensions.

\begin{table}[htbp]
\centering
\caption{Disease Classification Dimensions}
\label{tab:classification_dimensions}
\begin{tabular}{>{\bfseries}llp{8cm}}
\toprule
\textbf{Dimension} & \textbf{Values/Format} & \textbf{Description} \\
\midrule
Cause & \textit{Infectious, Genetic, Autoimmune, Neoplastic, Metabolic, Congenital, Toxic, Traumatic, Vascular, Degenerative, Nutritional, Psychiatric, Drug-induced, Iatrogenic, Idiopathic, Multifactorial} & Primary etiological category \\
Organ System & \textit{Heart, Lungs, Brain/Nervous, Liver, Kidneys, GI, Skin, Hematologic, Musculoskeletal, Eyes, Ears, Endocrine, Reproductive, Urinary, Immune/Lymphatic, Multiple/General} & Primary affected organ/system \\
Pathophysiology & Free text (standardized templates) & Mechanistic description of disease process \\
Duration & \textit{Acute, Subacute, Chronic, Episodic/Recurrent, Lifelong, Progressive, Variable} & Typical temporal pattern \\
Onset & \textit{Sudden, Insidious/Gradual, Congenital/Neonatal, Childhood/Adolescent, Adulthood, Pregnancy/Perinatal, Variable} & Typical onset pattern \\
Mode of Transmission & \textit{Person-to-person, Vector-borne, Environmental, Non-communicable (genetic, autoimmune, neoplastic, metabolic, traumatic, psychiatric, toxic), Food/water-borne, Blood-borne, Sexual, Vertical (mother-to-child)} & How disease spreads or is acquired \\
Clinical Course & \textit{Acute self-limited, Chronic progressive, Relapsing-remitting, Episodic/Recurrent, Lifelong stable, Progressive to terminal, Potentially curable, Variable} & Typical disease trajectory \\
Severity & \textit{Mild, Moderate, Severe/Life-threatening, Mild-to-Moderate, Moderate-to-Severe, Variable} & Typical severity range \\
Extent of Disease & \textit{Localized, Regional, Diffuse/Generalized, Metastatic/Disseminated, Variable} & Spatial distribution \\
Age Group & \textit{Neonates/Infants, Children/Adolescents, Young Adults, Adults, Elderly, All Ages, Pregnancy/Perinatal} & Typical affected age range \\
Epidemiology & \textit{Common, Rare (<1/2000), Age-related, Genetic/familial, Infectious (endemic/epidemic), Cancer (age-related), Autoimmune (female predominance), Variable} & Population-level patterns \\
ICD-11 Code & \textit{Valid ICD-11 code (e.g., 8A20, BA41.0, 5A11, 1C62)} & WHO ICD-11 classification code \\
\bottomrule
\end{tabular}
\end{table}

\section{Classification Coverage Statistics}
\label{sec:classification_stats}

\begin{table}[htbp]
\centering
\caption{Classification Coverage Across 841 Diseases}
\label{tab:coverage_stats}
\begin{tabular}{lrr}
\toprule
\textbf{Dimension} & \textbf{Diseases Classified} & \textbf{Coverage} \\
\midrule
Cause & 841 & 100\% \\
Organ System & 841 & 100\% \\
Pathophysiology & 841 & 100\% \\
Duration & 841 & 100\% \\
Onset & 841 & 100\% \\
Mode of Transmission & 841 & 100\% \\
Clinical Course & 841 & 100\% \\
Severity & 841 & 100\% \\
Extent of Disease & 841 & 100\% \\
Age Group & 841 & 100\% \\
Epidemiology & 841 & 100\% \\
ICD-11 Code & 841 & 100\% \\
\bottomrule
\end{tabular}
\end{table}

\section{Top Cause Categories}
\label{sec:cause_distribution}

\begin{table}[htbp]
\centering
\caption{Distribution of Primary Cause Categories}
\label{tab:cause_dist}
\begin{tabular}{lrr}
\toprule
\textbf{Cause Category} & \textbf{Count} & \textbf{Percentage} \\
\midrule
Infectious & 198 & 23.5\% \\
Neoplastic & 112 & 13.3\% \\
Autoimmune & 89 & 10.6\% \\
Genetic & 78 & 9.3\% \\
Metabolic/Endocrine & 72 & 8.6\% \\
Cardiovascular/Vascular & 61 & 7.3\% \\
Neurological/Degenerative & 54 & 6.4\% \\
Congenital & 47 & 5.6\% \\
Traumatic/Injury & 38 & 4.5\% \\
Toxic/Environmental & 34 & 4.0\% \\
Drug-induced/Iatrogenic & 28 & 3.3\% \\
Psychiatric & 22 & 2.6\% \\
Nutritional & 18 & 2.1\% \\
Idiopathic/Unknown & 15 & 1.8\% \\
Other & 15 & 1.8\% \\
\bottomrule
\end{tabular}
\end{table}

\section{Top Organ Systems}
\label{sec:organ_distribution}

\begin{table}[htbp]
\centering
\caption{Distribution of Primary Organ Systems}
\label{tab:organ_dist}
\begin{tabular}{lrr}
\toprule
\textbf{Organ System} & \textbf{Count} & \textbf{Percentage} \\
\midrule
Brain/Nervous System & 134 & 15.9\% \\
Heart/Cardiovascular & 98 & 11.7\% \\
Gastrointestinal & 87 & 10.3\% \\
Lungs/Respiratory & 72 & 8.6\% \\
Skin & 68 & 8.1\% \\
Hematologic/Immune & 61 & 7.3\% \\
Female Reproductive & 54 & 6.4\% \\
Kidneys/Urinary & 49 & 5.8\% \\
Endocrine & 44 & 5.2\% \\
Musculoskeletal & 41 & 4.9\% \\
Eyes & 28 & 3.3\% \\
Male Reproductive & 24 & 2.9\% \\
Ears & 18 & 2.1\% \\
Multiple/General & 36 & 4.3\% \\
\bottomrule
\end{tabular}
\end{table}

\section{Clinical Course Distribution}
\label{sec:course_distribution}

\begin{table}[htbp]
\centering
\caption{Clinical Course Patterns}
\label{tab:course_dist}
\begin{tabular}{lrr}
\toprule
\textbf{Course} & \textbf{Count} & \textbf{Percentage} \\
\midrule
Chronic, progressive/relapsing-remitting & 412 & 49.0\% \\
Acute, self-limited & 189 & 22.5\% \\
Episodic/Recurrent & 124 & 14.7\% \\
Variable & 78 & 9.3\% \\
Progressive to terminal & 24 & 2.9\% \\
Potentially curable & 15 & 1.8\% \\
\bottomrule
\end{tabular}
\end{table}

\section{Mode of Transmission Distribution}
\label{sec:transmission_distribution}

\begin{table}[htbp]
\centering
\caption{Mode of Transmission}
\label{tab:transmission_dist}
\begin{tabular}{lrr}
\toprule
\textbf{Transmission} & \textbf{Count} & \textbf{Percentage} \\
\midrule
Non-communicable & 689 & 81.9\% \\
Person-to-person / Vector-borne / Environmental & 152 & 18.1\% \\
\bottomrule
\end{tabular}
\end{table}

\section{Duration Distribution}
\label{sec:duration_distribution}

\begin{table}[htbp]
\centering
\caption{Disease Duration Patterns}
\label{tab:duration_dist}
\begin{tabular}{lrr}
\toprule
\textbf{Duration} & \textbf{Count} & \textbf{Percentage} \\
\midrule
Chronic (months to years) & 498 & 59.2\% \\
Acute (hours to days) & 214 & 25.4\% \\
Lifelong & 67 & 8.0\% \\
Episodic/Recurrent & 41 & 4.9\% \\
Progressive (years) & 15 & 1.8\% \\
Subacute (weeks to months) & 6 & 0.7\% \\
\bottomrule
\end{tabular}
\end{table}

\section{Onset Pattern Distribution}
\label{sec:onset_distribution}

\begin{table}[htbp]
\centering
\caption{Onset Patterns}
\label{tab:onset_dist}
\begin{tabular}{lrr}
\toprule
\textbf{Onset} & \textbf{Count} & \textbf{Percentage} \\
\midrule
Insidious/Gradual & 423 & 50.3\% \\
Sudden & 287 & 34.1\% \\
Adulthood & 64 & 7.6\% \\
Childhood/Adolescent & 38 & 4.5\% \\
Congenital/Neonatal & 29 & 3.4\% \\
Variable & 12 & 1.4\% \\
\bottomrule
\end{tabular}
\end{table}

\section{Severity Distribution}
\label{sec:severity_distribution}

\begin{table}[htbp]
\centering
\caption{Typical Severity Ranges}
\label{tab:severity_dist}
\begin{tabular}{lrr}
\toprule
\textbf{Severity} & \textbf{Count} & \textbf{Percentage} \\
\midrule
Variable & 434 & 51.6\% \\
Mild to Moderate & 156 & 18.5\% \\
Severe/Life-threatening & 124 & 14.7\% \\
Moderate to Severe & 78 & 9.3\% \\
Mild & 49 & 5.8\% \\
\bottomrule
\end{tabular}
\end{table}

\section{Extent of Disease Distribution}
\label{sec:extent_distribution}

\begin{table}[htbp]
\centering
\caption{Extent of Disease}
\label{tab:extent_dist}
\begin{tabular}{lrr}
\toprule
\textbf{Extent} & \textbf{Count} & \textbf{Percentage} \\
\midrule
Variable & 412 & 49.0\% \\
Localized & 287 & 34.1\% \\
Diffuse/Generalized & 109 & 13.0\% \\
Regional & 21 & 2.5\% \\
Metastatic/Disseminated & 12 & 1.4\% \\
\bottomrule
\end{tabular}
\end{table}

\section{Age Group Distribution}
\label{sec:age_distribution}

\begin{table}[htbp]
\centering
\caption{Primary Age Groups Affected}
\label{tab:age_dist}
\begin{tabular}{lrr}
\toprule
\textbf{Age Group} & \textbf{Count} & \textbf{Percentage} \\
\midrule
Adults & 398 & 47.3\% \\
Elderly & 187 & 22.2\% \\
All Ages & 112 & 13.3\% \\
Children/Adolescents & 89 & 10.6\% \\
Young Adults & 34 & 4.0\% \\
Neonates/Infants & 18 & 2.1\% \\
Pregnancy/Perinatal & 8 & 1.0\% \\
\bottomrule
\end{tabular}
\end{table}

\section{Sample Classification Entry}
\label{sec:sample_entry}

\begin{table}[htbp]
\centering
\caption{Example: Alzheimer Disease Classification}
\label{tab:alzheimer_example}
\begin{tabular}{>{\bfseries}ll}
\toprule
\textbf{Dimension} & \textbf{Value} \\
\midrule
Cause & Genetic \\
Organ System & Brain \\
Pathophysiology & Protein aggregation $\to$ synaptic loss $\to$ neuronal death $\to$ atrophy \\
Duration & Chronic progressive (8-10 years avg) \\
Onset & Insidious ($>$65 late-onset; $<$65 early-onset) \\
Mode of Transmission & Non-communicable \\
Clinical Course & Chronic, progressive/relapsing-remitting \\
Severity & Mild to severe (total dependence) \\
Extent of Disease & Diffuse cortical (temporal, parietal, frontal) \\
Age Group & $>$65 years (95\%); early-onset $<$65 \\
Epidemiology & 6.7M US (2023); 55M global; doubling every 20 years \\
ICD-11 Code & 8A20 \\
\bottomrule
\end{tabular}
\end{table}

\section{Implementation Notes}
\label{sec:impl_notes}

\begin{itemize}[noitemsep]
\item All 841 disease chapters include a \texttt{Classification} section immediately after \texttt{Overview}.
\item Classification uses a 12-field description list (\texttt{description} environment).
\item ICD-11 codes mapped from WHO ICD-11 MMS (Mortality and Morbidity Statistics).
\item Cause categories derived from WHO Global Health Estimates and ICD-11 chapter structure.
\item Organ system mapping follows ICD-11 body system chapters.
\item Epidemiology data sourced from WHO GHE 2019, CDC, NIH, and disease-specific registries.
\item ICD-11 codes validated against \url{https://icd.who.int/} (2024 release).
\end{itemize}